% Лекц 15: семинар ба гэрийн даалгавар. Бодлогууд: Г.Н. Берман, "Сборник задач по курсу
% математического анализа" (22-р хэвлэл). [Б.N] нь номын бодлогын дугаар.

\ProblemSet{Семинар 15}
\begin{frame}[allowframebreaks]{Семинар 15}
\pgroup{Интеграл нийлбэр, тодорхой интегралын чанар.}
\begin{plist}
\prob{1616}{$\int_0^1e^x\,dx$ интегралыг шууд нийлбэрийн хязгаараар бод (хэрчмийг $n$ тэнцүү хэсэгт хуваа).}
\prob{1621}{$\int_1^2\frac{dx}{x}$-ийн хувьд хэрчмийг $n$ тэнцүү хэсэгт хувааж интеграл нийлбэр зохиогоод $\displaystyle\lim_{n\to\infty}\Bigl(\frac1n+\frac{1}{n+1}+\dots+\frac{1}{2n}\Bigr)$-г ол.}
\prob{1628}{$\int_0^{10}\frac{x\,dx}{x^3+16}<\frac56$ болохыг батал.}
\prob{1636}{Бодолгүйгээр аль интеграл их болохыг тогтоо: 1) $\int_0^1x^2dx$ ба $\int_0^1x^3dx$; \ 2) $\int_1^2x^2dx$ ба $\int_1^2x^3dx$.}
\prob{1644}{$y=2x^2+3x+3$ функцийн $[1,4]$ хэрчим дээрх дундаж утгыг ол.}
\end{plist}
\pgroup{Хувьсах дээд хилтэй интеграл.}
\begin{plist}
\prob{1658}{$y=\int_0^x\frac{1-t+t^2}{1+t+t^2}\,dt$ функцийн $x=1$ дээрх уламжлалыг ол.}
\prob{1662}{$y=\int_0^{2x}\frac{\sin t}{t}\,dt$ функцийн $x$-ээр авсан уламжлалыг ол.}
\prob{1663.1}{$\frac{d}{dx}\int_2^{e^x}\frac{\ln z}{z}\,dz$-г ол.}
\prob{1668}{$x$-ийн ямар утгад $I(x)=\int_0^xte^{-t^2}dt$ функц экстремумтай вэ? Түүний утга?}
\end{plist}
\pgroup{Ньютон–Лейбницийн томьёогоор бод.}
\begin{plist}
\prob{1672.4}{$\int_1^2\bigl(x+\frac1x\bigr)^2dx$}
\prob{1673.5}{$\int_0^1\frac{dx}{1+x^2}$}
\prob{2231}{$\int_0^1\sqrt{1+x}\,dx$}
\prob{2237}{$\int_0^1(e^x-1)^4e^x\,dx$}
\prob{2239}{$\int_0^1\frac{x\,dx}{(x^2+1)^2}$}
\prob{2248}{$\int_0^1\frac{dx}{x^2+4x+5}$}
\end{plist}
\pgroup{Хэсэгчлэн интегралчлах ба орлуулах арга.}
\begin{plist}
\prob{2259}{$\int_0^1xe^{-x}dx$}
\prob{2260}{$\int_0^{\pi/2}x\cos x\,dx$}
\prob{2264}{$\int_0^{e-1}\ln(x+1)\,dx$}
\prob{2275}{$\int_4^9\frac{\sqrt x}{\sqrt x-1}\,dx$}
\prob{2276}{$\int_0^1\frac{\sqrt x\,dx}{1+x}$}
\end{plist}
\end{frame}

\ProblemSet{Гэрийн даалгавар 15}
\begin{frame}[allowframebreaks]{Гэрийн даалгавар 15}
\pgroup{Талбай, үнэлгээ, дундаж утга.}
\begin{plist}
\prob{1625}{$y=x^3$ ба $y=x$ шугамуудаар хүрээлэгдсэн дүрсийн талбайг ол.}
\prob{1626}{$y=x^2-2x-3$ ба $y=-x^2+6x-3$ параболуудаар хүрээлэгдсэн дүрсийн талбайг ол.}
\prob{1630}{$\int_{1{,}5}^{3{,}5}\frac{x^2\,dx}{x-1}$ интегралыг үнэл.}
\prob{1637}{Аль интеграл их вэ: 1) $\int_0^12^{x^2}dx$ ба $\int_0^12^{x^3}dx$; \ 2) $\int_1^2\ln x\,dx$ ба $\int_1^2(\ln x)^2dx$?}
\prob{1642}{$y=kx+b$ функцийн $[x_1,x_2]$ дээрх дундаж утгыг олж, функц уг утгыг авах цэгийг ол.}
\prob{1645}{Геометрийн үүднээс $y=\sqrt{a^2-x^2}$ функцийн $[-a,a]$ дээрх дундаж утгыг ол.}
\end{plist}
\pgroup{Хувьсах хилтэй интеграл.}
\begin{plist}
\prob{1650.2}{$\int_a^xt^5\,dt$-г бод.}
\prob{1659}{$y=\int_0^x\sin t\,dt$ функцийн $x=0,\ \frac{\pi}{4},\ \frac{\pi}{2}$ цэгүүд дээрх уламжлалыг ол.}
\prob{1661}{$y=\int_x^5\sqrt{1+t^2}\,dt$ функцийн $x=0$ ба $x=\frac34$ дээрх уламжлалыг ол.}
\prob{1663.2}{$\frac{d}{dx}\int_{x^2}^1\ln t\,dt$-г ол.}
\prob{1667}{$y=\int_0^{z^3}\frac{dx}{1+x^3}$ функцийн $z$-ээр авсан хоёрдугаар эрэмбийн уламжлалыг $z=1$ дээр ол.}
\prob{1670}{$y=\int_0^x(t^2-3t+2)\,dt$ функцийн экстремум ба нугаралтын цэгийг олж, графикийг байгуул.}
\end{plist}
\pgroup{Интегралыг бод.}
\begin{plist}
\prob{1672.6}{$\int_1^2(\sqrt x-\sqrt[3]{x})\,dx$}
\prob{1673.4}{$\int_0^{\pi/4}\frac{dx}{\cos^2x}$}
\prob{2232}{$\int_{-2}^{-1}\frac{dx}{(11+5x)^3}$}
\prob{2234}{$\int_4^9\frac{y-1}{\sqrt y+1}\,dy$}
\prob{2240}{$\int_1^e\frac{dx}{x\sqrt{1-(\ln x)^2}}$}
\prob{2242}{$\int_1^2\frac{e^{1/x}}{x^2}\,dx$}
\prob{2247}{$\int_2^3\frac{dx}{2x^2+3x-2}$}
\prob{2251}{$\int_{-\pi/2}^{\pi/2}\frac{dx}{1+\cos x}$}
\prob{2252}{$\int_0^{\pi/2}\cos^5x\sin2x\,dx$}
\prob{2261}{$\int_{\pi/4}^{\pi/3}\frac{x\,dx}{\sin^2x}$}
\prob{2263}{$\int_1^2x\log_2x\,dx$}
\prob{2267}{$\int_0^{\pi/2}e^{2x}\cos x\,dx$}
\prob{2268}{$\int_1^e\ln^3x\,dx$}
\prob{2277}{$\int_3^8\frac{x\,dx}{\sqrt{1+x}}$}
\prob{2278}{$\int_0^1\frac{x\,dx}{1+\sqrt x}$}
\prob{2286}{$\int_1^2\frac{\sqrt{x^2-1}}{x}\,dx$}
\prob{2288}{$\int_0^1\sqrt{(1-x^2)^3}\,dx$}
\prob{2296}{$y=\sqrt x+\frac{1}{\sqrt x}$ функцийн $[1,4]$ дээрх дундаж утгыг ол.}
\end{plist}
\end{frame}
